% ==================================================================
% id:       comb-022
% topic:    קומבינטוריקה
% subtopic: המקדם הבינומי
% solution: full
% ==================================================================
\begin{enumerate}
\item הראו שמספר המילים באורך $n$ מעל $\set{0,1}$ שיש בהן בדיוק $r$ אפסים הוא $\binom nr$.
\begin{solbox}
מילה כזו נקבעת לגמרי לפי המקומות שבהם יש $0$ (בכל שאר המקומות יש $1$). לכן צריך לבחור $r$ מקומות מתוך $n$, בלי חשיבות לסדר ובלי חזרות, ויש לכך $\binom nr$ אפשרויות.
\end{solbox}
\item כמה סדרות של אפסים ואחדים באורך $2n$ יש, שבהן מספר האפסים ב-$n$ המקומות הראשונים שווה למספר האחדים ב-$n$ המקומות האחרונים?
\hint{אם ב-$n$ המקומות הראשונים יש $i$ אפסים, כמה אפסים יש בסדרה כולה?}
\begin{solbox}
נניח שב-$n$ המקומות הראשונים יש $i$ אפסים. לפי התנאי, ב-$n$ המקומות האחרונים יש $i$ אחדים, כלומר $n-i$ אפסים. לכן בסדרה כולה יש $i+(n-i)=n$ אפסים.
להפך, אם בסדרה יש בדיוק $n$ אפסים, ובחצי הראשון יש $i$ מהם, אז בחצי השני יש $n-i$ אפסים, כלומר $i$ אחדים, והתנאי מתקיים.
לכן אלה בדיוק הסדרות באורך $2n$ עם $n$ אפסים, ולפי הסעיף הקודם יש $\ans{\binom{2n}n}$ כאלה.
\end{solbox}
\end{enumerate}
